%---------------------------------------------------------------------------%
%->> Main content
%---------------------------------------------------------------------------%

\section{选题的背景及意义}

\subsection{长上下文推理的服务需求}

近年来，大语言模型从通用对话逐步进入长文档理解、代码仓库分析、检索增强生成和智能体任务规划等场景。与普通问答不同，这类任务需要模型一次性读入大量历史信息、外部知识或工具调用轨迹，在线请求的上下文长度从 4K、8K 扩展到 128K、256K 乃至百万 token 级别。LLaMA、Qwen、DeepSeek 等开源模型持续扩大上下文窗口和模型规模\citep{touvron2023llama,touvron2023llama2,dubey2024llama3,yang2024qwen2,yang2025qwen25,deepseekai2024deepseekv2,deepseekai2025deepseekv3,deepseekai2025deepseekr1}，长上下文能力已从评测指标演变为推理服务必须稳定支撑的负载类型。

长上下文推理同时改变了单请求成本和在线负载结构。长请求推高首 token 延迟（Time To First Token, TTFT），并在生成阶段持续增加每个输出 token 的缓存访问量；但真实服务中并非只有长请求——短对话、RAG 查询、长文档分析和智能体任务通常混合到达。系统若只按请求数量分配负载，会把一个超长请求和一个短请求视作近似相同的工作量，进而在数据并行副本之间产生实际执行时间差。本文要解决的核心问题由此产生：面对长短请求混合到达的负载，如何让长请求获得必要的并行度分摊访存压力，同时不让短请求被额外通信开销拖慢。随着上下文长度继续扩大，这个问题的瓶颈不取决于单点算子效率，而由并行策略、缓存管理、调度机制和硬件带宽共同决定。

\subsection{长上下文推理的资源瓶颈}

自回归大模型推理分为 Prefill 和 Decode 两个阶段。Prefill 阶段一次性处理输入 token，矩阵乘规模大，主要表现为计算密集，影响长上下文的 TTFT；Decode 阶段逐 token 生成，每一步读取模型权重和历史 KV Cache，计算强度低，受 HBM 带宽约束，直接影响 Time Per Output Token（TPOT）。对于 $n_\text{layers}$ 层模型，每层 $n_\text{kv\_heads}$ 个 KV head、head 维度 $d_\text{head}$、数据类型 $B_\text{dtype}$ 字节，长度 $L$ 的请求所需 KV Cache 容量为：
\[
F(L)=2 \times n_\text{layers} \times L \times n_\text{kv\_heads} \times d_\text{head} \times B_\text{dtype}.
\]
该式直接刻画了 Decode 阶段的访存增长规律：单步 KV Cache 读取量随 $L$ 线性增大。以 128K 上下文对比 4K 为例，仅 KV Cache 读取量即相差约 32 倍。当长短请求混合到达时，仅按请求数量进行数据并行（Data Parallelism, DP）路由掩盖了真实负载差异——200K token 请求的 KV Cache 占用和 HBM 读取时间远高于普通短请求，即使各 DP rank 承载相同请求数，其 KV Cache 占用、HBM 访问量和执行时间仍可能悬殊。

现有推理系统通过 continuous batching、PagedAttention、chunked prefill、FlashAttention/FlashDecoding 等方法提高吞吐或降低显存碎片\citep{kwon2023efficient,dao2022flashattention,dao2023flashattention2,flashdecoding2023}，中短上下文场景效果显著。但超长上下文下仍面临两个限制：KV Cache 总量随序列长度增长的趋势未变；长短请求混合后，瓶颈从单请求访存扩展为多副本间的负载不均。长上下文推理需要在缓存组织和请求调度层面进一步优化。

\subsection{序列并行与静态配置问题}

序列并行（Context Parallelism, CP）沿序列维度切分输入或 KV Cache，直接缓解长序列的计算和 KV Cache 容量瓶颈。PCP（Prefill Context Parallelism）将长 prompt 的 Prefill 计算分布到多个 rank；DCP（Decode Context Parallelism）将 Decode 阶段的 KV Cache 按序列维度分片，每个 rank 只读取一部分历史 KV，降低单 rank 的显存容量和 HBM 带宽压力\citep{liu2023ringattention,jacobs2023deepspeedulysses,fang2024usp,zhong2024contextparallelism,vllmascendcp2026}。但现有 CP 在服务系统中多为静态配置：服务启动前设定统一 CP size，所有请求共享同一配置。这与混合负载需求直接冲突——短请求本可由单个 DP rank 低延迟完成，被迫走 CP 路径会承担不必要的跨 rank 通信；长请求若不走 CP，则面临 KV Cache 容量不足和 HBM 访存过重。

动态序列并行是本文的主要技术切入点，思路是在运行时根据上下文长度为每个请求独立选择 CP/DP 执行路径：短请求保持 CP=1 的单 DP 路径，避免跨 rank 通信；长请求触发 CP$>$1，由多个 DP rank 协同分摊 KV Cache、attention 计算和访存压力。前期工作在 vLLM 社区提出了 Dynamic Context Parallelism RFC\citep{vllmdynamiccp2025}，阐述了 Domain、跨 DP 调度和 attention 聚合等设计概念；本文在此基础上完成系统实现、昇腾平台适配和实验验证。

\subsection{MoE 同步等待对负载不均的放大效应}

混合专家模型（Mixture of Experts, MoE）将 rank 间执行时间差进一步放大为同步等待。DeepSeek-V2、DeepSeek-V3、DeepSeek-R1 等 MoE 模型通过稀疏专家层扩大模型容量，MoE 层依赖 all-to-all 通信完成 token 到专家的分发和聚合\citep{lepikhin2020gshard,fedus2021switch,deepseekai2024deepseekv2,deepseekai2025deepseekv3,deepseekai2025deepseekr1}。all-to-all 形成同步屏障：同一同步组内较快的 DP rank 必须等待最慢 rank 到达 MoE 层后才能继续执行。当某个 rank 因长上下文请求在 attention 阶段遍历大量 KV Cache，其他 rank 在 MoE 同步点空转，形成 MoE bubble（如图~\ref{fig:moe-long-request-blocking}），降低集群吞吐并恶化尾延迟。静态 CP 无法消除这种不均衡：CP size 统一为 1 时，长请求仍由单 rank 独自承担全部 KV 访存；统一启用 CP 时，短请求又因不必要的跨 rank 通信而延迟增大。缓解 MoE bubble 需要请求级 CP/DP 路径选择与跨 rank 协同调度配合。

\begin{figure}[!htbp]
    \centering
    \includegraphics[width=\textwidth]{1-moe-bubble}
    \caption{长请求导致的 MoE 同步等待}
    \label{fig:moe-long-request-blocking}
\end{figure}
\FloatBarrier

\subsection{昇腾平台的硬件特征与适配需求}

昇腾平台的长上下文推理还需要处理硬件和软件栈的适配差异。昇腾生态中的 CANN 图执行、HCCL/HCCS 通信、torch-npu 算子接口、vLLM-Ascend 运行时和 KV Cache 布局，与 CUDA 生态的常见实现不完全一致——通信域组织、图捕获参数形状、attention 后端输入布局和缓存地址映射等 GPU 上可直接复用的假设，在昇腾上均需重新检查。

昇腾的硬件特征同时也为动态 CP 提供了有利条件。以实验所用的 Atlas A3（Ascend 910C）为例，其 FP16 峰值算力与 HBM 带宽之比约为同代 NVIDIA GPU 的 1.5--2 倍，即同等计算能力下 HBM 带宽更低，长上下文 Decode 的访存瓶颈比 GPU 更突出。动态 CP 通过多 rank 分摊 KV Cache 访存来降低单 rank 带宽压力，在昇腾上的边际收益更大，启用 CP 的上下文长度阈值可相应设得更低。此外，昇腾超节点架构中节点间互联带宽与节点内 HCCS 带宽一致，跨节点通信不会出现 NVIDIA 平台 NVLink 到 InfiniBand 的带宽断崖，CP rank group 可跨节点组建以支持更大 CP size，覆盖更长上下文场景。

上述差异在 Decode 阶段尤为明显。长上下文 Decode 对 HBM 带宽和 KV Cache 访问路径高度敏感，若调度器把长请求集中到少数 rank，这些 rank 在 attention 阶段花费更多时间读取历史 KV，即便其他 rank 已完成本地计算，也会被后续 MoE 同步通信拖住。昇腾平台上的优化不能只停留在算子替换或接口适配，需要将缓存管理、跨 rank 调度和通信范围控制纳入统一的系统设计。

\subsection{研究问题}

综上，本文的核心问题是：混合负载下如何做到长请求获得必要的序列维度并行，同时不让短请求承担额外通信开销。这个问题分解为三个相互依赖的子问题。

\textbf{请求级 CP/DP 路径判定与跨 DP 协同调度。}动态 CP 请求必须在所有参与 rank 的同一调度周期中同步出现，否则导致通信死锁或结果错误；传统各 rank 本地独立编排批次的 scheduler 缺乏这种跨 rank 全局协调能力。静态 CP 将服务级配置强加于所有请求，无法兼顾短请求低延迟和长请求容量分摊。路径判定与调度因此需要整体设计——中心化调度器按上下文长度选择 CP/DP 路径，结合各 rank 剩余容量和 KV block 分配参与 rank 集合，生成按 rank 组织的调度输出。Domain 执行架构作为调度与执行间的桥接层，负责输出分发和跨 rank 元数据同步。此外，现有 DCP 将 TP 与 CP 绑定，单个 rank 无法独立执行不走 CP 的请求，需要解耦 Decode 阶段的 TP 与 CP，使 CP 组在 TP 组之上按需动态组建。

\textbf{跨 DP KV Cache 管理。}动态 CP 使同一请求的 KV Cache 跨多个 DP rank 物理分片存储，而上层 attention 后端仍需按逻辑连续的 KV 序列访问——这一"物理分布、逻辑连续"的矛盾是现有单 rank KV Cache 管理方案未涉及的，也是研究内容二要统一解决的核心问题。该矛盾在三个子场景中分别体现：（1）缓存寻址——需通过二维 slot mapping 和 virtual block 将多层物理 block 组织为统一逻辑 KV 空间，支持 CP 与 DP 请求共存；（2）去冗余——原 DCP 假设所有请求均走 CP 路径以执行 TP 域 KV 去冗余，DyCP 引入 DP 请求后需按请求类型区分 KV 布局；（3）跨节点传输——PD 分离下长请求 KV 按 CP 分片跨 rank 分布，需按分片粒度传输并在 Decode 端重建虚拟 KV 视图。三个子场景回扣同一矛盾：物理层 KV 被 CP 分片打散后，逻辑层如何维持连续一致的访问视图。

\textbf{CP+DP 混合执行。}同一 Decode 批次可能同时包含 CP 请求和 DP 请求。若所有请求参与跨 rank attention 通信，短请求被通信拖慢；若通信区域随 CP/DP 比例变化，图捕获参数形状不稳定导致编译缓存命中率下降。混合执行需要解决三个适配问题：attention 后端需针对 MLA latent KV 分片粒度适配 CP 索引与 LSE 聚合（首先适配 MLA，后续扩展至 GQA/MHA）；chunked prefill 下 Prefill/Decode $\times$ CP/DP 形成 2$\times$2 执行矩阵，需要元数据拆分与输出拼接；批次重排（CP 前置、DP 后置）与二维 AclGraph 捕获需协同设计，限定通信区域同时控制图编译开销。

三个问题构成从判定到执行的完整链路：研究内容一决定何时启用 CP 并保证多 rank 协同；研究内容二在"物理分布、逻辑连续"的约束下统一管理跨 DP KV Cache；研究内容三在混合批次执行中实现 CP/DP 通信隔离与正确聚合。

\subsection{选题意义}

在服务效率方面，在线推理服务由短对话、RAG 查询、长文档分析和智能体任务混合构成，静态 CP 难以适配这种负载差异：长请求需要多 rank 分摊 KV Cache 和访存压力，短请求更适合留在单 rank 本地执行。动态 CP 通过请求级路由——短请求不通信、长请求才通信——避免短请求尾延迟被 CP 通信放大，同时缓解长请求的容量和访存压力。

在平台层面，本文面向昇腾平台展开，有助于补齐国产 AI 加速器上长上下文推理服务的工程经验。昇腾超节点架构中节点间与节点内通信带宽一致，CP rank group 可跨节点组建以覆盖超长上下文；昇腾 NPU 的计算访存比约为同代 GPU 的 1.5--2 倍，访存瓶颈更突出，动态 CP 分摊 KV 访存的收益边界更宽、启用阈值更低。均匀带宽降低了跨节点 CP 通信代价，高计算访存比扩大了 CP 分摊访存的收益区间，两者使昇腾平台适合验证动态 CP 的实际效果。

从研究内容看，本文将序列并行从"服务级静态配置"推进到"请求级动态调度"，贯通调度器、KV Cache 管理、通信范围控制和执行后端四个层面，为昇腾平台上的长上下文推理优化提供可复用的系统设计。

\section{国内外本学科领域的发展现状与趋势}

\subsection{长上下文推理服务优化研究}

大模型推理服务系统的基础优化主要包括 continuous batching、PagedAttention、chunked prefill、FlashAttention/FlashDecoding、算子融合和显存复用等。PagedAttention 借鉴操作系统分页思想，将 KV Cache 划分为固定大小 block，通过 block table 管理逻辑序列到物理缓存的映射，减少显存碎片并提升多请求服务能力\citep{kwon2023efficient}。FlashAttention 系列从访存路径和并行划分角度优化 attention 计算\citep{dao2022flashattention,dao2023flashattention2,flashdecoding2023}。SGLang、vLLM、DeepSpeed Inference、FlexGen 等系统则分别从请求调度、结构化执行、多 GPU 推理和异构内存 offloading 等方向提升服务吞吐\citep{aminabadi2022deepspeedinference,sheng2023flexgen,zheng2023sglang}。

随着上下文长度扩大，KV Cache 已成为制约推理服务吞吐的关键瓶颈。MQA、GQA 和 MLA 等注意力结构降低了每 token 的 KV Cache 开销\citep{shazeer2019fast,ainslie2023gqa,deepseekai2024deepseekv2}。H2O、Scissorhands、KIVI、KVQuant 和 InfiniGen 等工作从 token 保留、量化和预取角度减少 KV Cache 容量或访存压力\citep{zhang2023h2o,liu2023scissorhands,liu2024kivi,hooper2024kvquant,lee2024infinigen}。然而，上述方法聚焦于降低单请求或单步的访存开销，并未考虑混合负载中长短请求在不同 DP rank 上的调度差异。在 MoE 服务场景下，即使 KV Cache 容量和访存已优化，若不同 rank 的请求长度差异过大，仍会在 all-to-all 同步点形成等待，所以 KV Cache 优化还需要与跨 rank 调度协同设计。

\subsection{阶段异构与推理调度研究现状}

Splitwise、DistServe 和 Mooncake 等工作表明，长上下文推理不能只依赖单点算子优化，服务系统还要处理 Prefill/Decode 阶段、请求队列和缓存状态之间的耦合，其共同趋势是将推理服务从"单实例内尽量合批"推进到"按阶段和资源特征组织服务"。但这些工作主要解决 Prefill/Decode 阶段间的资源错配或 KV Cache 生命周期管理，并不回答同一在线服务内长短请求如何混合调度。对于 MoE Decode 批次，多个 DP rank 仍会因上下文长度不同产生 KV 使用量和 attention 完成时间差，进而转化为 all-to-all 同步等待。因此，本文聚焦请求级 CP/DP 路径选择和跨 DP 协同调度。

\subsection{序列并行研究现状}

序列并行沿序列维度切分长上下文，与按样本划分的数据并行、按矩阵维度划分的张量并行和按层划分的流水并行不同，CP 直接针对长序列本身的计算和 KV Cache 容量瓶颈。Ring Attention 通过环形传递 KV block 支持超长序列训练和推理\citep{liu2023ringattention}；DeepSpeed-Ulysses 和 USP 等工作从序列并行和多维并行角度提高长序列可扩展性\citep{jacobs2023deepspeedulysses,fang2024usp}；Context Parallelism for Scalable Million-Token Inference 等工作表明 CP 适用于百万 token 推理场景\citep{zhong2024contextparallelism}。vLLM-Ascend 文档展示了 PCP/DCP 在工程框架中的 rank layout、slot mapping、virtual block 和通信组设计\citep{vllmascendcp2026}。

现有 CP 的主要局限在于配置静态、执行路径固定：服务启动时设定的 CP size 对所有请求统一生效，无法按单请求长度做 CP=1 与 CP$>$1 的运行时切换。部分 DCP 实现还与 TP 强绑定，每个 rank 同时承担序列切分和张量切分职责，单个 rank 无法独立执行不需 CP 的短请求，CP 请求和 DP 请求难以在同一批次中混合执行，只能采用"全批次 CP"或"全批次 DP"的粗粒度策略。Dynamic Context Parallelism RFC 是作者前期在 vLLM 社区提出的设计方案\citep{vllmdynamiccp2025}，思路是请求级异构并行——短请求走常规 DP 路径，长请求由多个 DP attention 副本协同执行。本文在此方案基础上进一步完成中心化调度、跨 DP KV 管理、混合批次图执行和昇腾平台适配等系统工程问题。

\subsection{MoE 推理的负载均衡与同步优化}

MoE 模型在训练中常通过 capacity factor、辅助负载均衡损失和 expert placement 减少专家负载不均\citep{lepikhin2020gshard,fedus2021switch}。推理阶段的难点不同：路由分布由输入内容决定，Prefill 阶段长请求的计算量远大于短请求，Decode 阶段长请求的 KV Cache 访存量远大于短请求，两个阶段都会在 DP rank 之间产生执行时间差。MoE all-to-all 会把这种差异放大为同步等待。已有优化主要围绕专家放置、路由和 all-to-all 通信协议本身，对 attention/KV 侧负载差异导致的 rank 间到达时间差关注较少。

动态 CP 从 attention/KV 侧处理这一等待问题：长请求不再由单一 DP rank 独自遍历完整 KV Cache，而是由多个 rank 沿序列维度分摊 attention 和 KV 访存，使参与 rank 以更接近的时间到达 MoE 同步点。CP 在扩展容量之外，同时降低 MoE 同步等待，与专家侧负载均衡形成互补。这也意味着动态 CP 的改动不局限于 attention 算子，而需协同处理路径判定、KV Cache 跨 rank 分布和通信范围控制。

综上，现有工作在 KV Cache 压缩、阶段异构调度、序列并行和 MoE 通信优化四个方向上已取得显著进展，但缺少一个将序列并行从静态配置推进为请求级动态调度的系统方案。这正是本文的研究切入点：在长短混合负载下，通过动态 CP 让长请求获得必要并行度、短请求不受额外通信影响，并将判定、管理、执行三个环节作为统一系统进行设计。以下 §3 对此展开。

\section{课题主要研究内容、预期目标}

\subsection{主要研究内容}

本文围绕动态序列并行这一核心机制，从"判定→管理→执行"链路展开以下三方面研究。

\textbf{研究内容一：请求级动态 CP 判定与中心化跨 DP 协同调度。}现有序列并行方案的 CP size 在服务启动时设定，运行期间所有请求共享同一配置，无法兼顾短请求低延迟和长请求容量分摊；动态 CP 请求必须在所有参与 rank 的同一调度周期同步出现，而传统各 rank 本地 scheduler 缺乏跨 rank 全局协调能力。研究内容包括：基于上下文长度阈值与硬件计算访存比的请求级 CP/DP 路径判定方法；中心化跨 DP 协同调度机制，统一管理全局请求队列与各 rank 资源状态，生成按 rank 组织的调度输出并保证 CP 请求的调度一致性；Domain 执行架构（DomainEngineCoreProc）作为调度与多进程执行的桥接层，负责输出分发与跨 rank 元数据同步；Decode 阶段 TP 与 CP 的解耦，CP 组在 TP 组之上按需动态组建。研究重点包括阈值与平台特征的定量关系、长短请求分离队列设计、CP rank 集合选择与通信拓扑的结合、TP-CP 解耦后通信组的动态管理，以及该调度方式对 MoE all-to-all 同步等待的缓解效果。

\begin{figure}[htbp]
    \centering
    \includegraphics[width=\textwidth]{corssdpscheduler}
    \caption{CrossDPScheduler 的跨 DP 调度架构}
    \label{fig:crossdpscheduler-overview}
\end{figure}

\textbf{研究内容二：支持动态并行粒度的跨 DP KV Cache 管理。}动态 CP 使同一请求的 KV Cache 跨多个 DP rank 物理分片存储，而上层 attention 后端仍需按逻辑连续的 KV 序列访问。这一"物理分布、逻辑连续"的矛盾在三个子场景中分别展开，是本项研究内容要统一解决的问题。

（1）跨 DP KV Cache 虚拟化与混合共存。CP 请求的 KV 分布在多个 rank 的本地 block pool 中，DP 请求的 KV 仅驻留在单个 rank。二维 slot mapping（第一维为参与 rank 编号，第二维为逻辑 block 区间）将多个 rank 的物理 KV block 组织为统一的逻辑 KV 空间，virtual block 在逻辑层提供跨 rank 连续视图，mask 标记区分 CP 与 DP 请求的缓存区域。方案保留各 rank 独立 block pool 以兼容 vLLM 现有 block 级内存管理，virtual block 仅在逻辑层引入映射而不侵入底层分配器。

（2）TP 域内 KV Cache 去冗余的分类型适配。原 DCP 假设同一 TP 组内所有请求均走 CP 路径，KV head 相同时仅保留一份副本。DyCP 引入 DP 请求后，CP 请求的 KV 已按序列维度跨 DP rank 分片（每 rank 仅持部分），DP 请求的 KV 完整保留在本地。去冗余逻辑需按请求类型区分：DP 请求沿用原策略，CP 请求在跨 DP 分片后的局部 KV 上执行去冗余，防止误删 DP 请求的完整 KV 或遗漏 CP 请求分片 KV 的冗余副本。

（3）PD 分离下的分片级 KV 传输。Prefill 产出的长请求 KV 按 CP 分片分布在多个 rank，与传统 PD 分离"每个请求 KV 完整存储在单 rank"的假设不同。在已有 Mooncake connector 基础上扩展按分片粒度的传输：每个 Prefill rank 仅传输本地持有的 KV 分片至 Decode 端对应 rank，Decode 端接收后通过 slot mapping 重建虚拟 KV 视图，复用已有连接管理和序列化机制。

三个子场景回扣同一核心矛盾：物理层因 CP 分片而跨 rank 分布后，逻辑层如何维持连续一致的 KV 视图。研究重点包括 slot mapping 的正确性和一致性、virtual block 与已有 block table 的兼容性、TP 去冗余与跨 DP 分片的不冲突性，以及分片传输的 KV 重建开销。

\textbf{研究内容三：CP+DP 混合执行。}动态 CP 引入 CP 请求与 DP 请求在同一批次中共存的执行模式，需在 attention 计算、图捕获和批次管理三方面适配。CP attention 后端首先适配 MLA 结构（mla\_cp），处理 latent KV 分片索引与跨 rank LSE 聚合，并在 chunked prefill 下支持 Prefill/Decode $\times$ CP/DP 的 2$\times$2 执行矩阵的元数据拆分与输出拼接，后续扩展至 GQA、MHA 等结构。批次管理采用 CP 请求前置、DP 请求后置的重排策略限定通信区域，聚合完成后通过 hidden state restore 恢复原始请求顺序，保证后续 FFN/MoE 层输入正确。图捕获方面，CP+DP 混合批次使图参数形状由 batch token 数和 CP 请求数两个维度共同决定，需设计二维 bucketing 策略在图缓存大小与命中率之间平衡。研究重点包括 CP attention 后端的 KV 分片与 LSE 聚合正确性、2$\times$2 执行矩阵下的输出正确性、批次重排对短请求尾延迟的影响，以及二维图缓存的编译开销控制。

\begin{figure}[htbp]
    \centering
    \includegraphics[width=0.92\textwidth]{batchreorder}
    \caption{CP+DP 混合批次的 batch reorder 与 hidden state restore}
    \label{fig:batchreorder-overview}
\end{figure}

本文的主要研究贡献包括：
\begin{enumerate}
    \item 在 vLLM-Ascend 调度层实现请求级 CP/DP 判定与中心化跨 DP 调度器（CrossDPScheduler），设计 Domain 执行架构（DomainEngineCoreProc）统一管理调度输出分发与跨 rank 元数据同步，并解耦 Decode 阶段 TP 与 CP 使同一组 Worker 可灵活切换 CP/DP 路径。
    \item 扩展 vLLM 的 KV Cache 管理层，实现跨 DP KV Cache 虚拟化管理器（CrossDPKVCacheManager），以"物理分布、逻辑连续"为统一框架解决 KV 虚拟化、TP 去冗余适配和 PD 分离分片传输三个子问题。
    \item 在昇腾平台实现 CP attention 后端，首先适配 MLA 结构的 latent KV 分片与跨 rank 聚合，后续扩展至其他注意力结构；在 AclGraph 图模式下实现 CP+DP 混合批次的二维图捕获、批次重排与 hidden state restore。
\end{enumerate}

\subsection{预期目标}

本文预期达成以下研究目标：

\textbf{系统原型交付。}在昇腾推理框架（vLLM-Ascend）中实现动态 CP 原型系统，包含请求级 CP/DP 判定模块、中心化跨 DP 调度器与 Domain 执行架构、跨 DP KV Cache 虚拟化管理器（含 TP 域 KV 去冗余适配和 PD 分离 KV 传输适配）、CP attention 后端（首先适配 MLA 结构）和 CP+DP 混合批次执行引擎，形成从请求到达到推理输出的完整动态 CP 执行链路。

\textbf{服务性能提升。}前期在 16 卡 Ascend 910C 上以 DeepSeek-V2-Lite 进行的 DP 与静态 DCP 基线实验（图~\ref{fig:dp-dcp-long-context}）已显示，静态 DCP 在低负载长上下文下可降低 P90 尾延迟，但在高负载下因所有请求均参与通信而劣于纯 DP。动态 CP 的目标是在混合负载下同时改善两端：相比纯 DP 降低长请求的 MoE 同步等待时间和 TPOT 尾延迟，相比静态 CP 避免高负载下短请求的通信惩罚。通过消融实验量化各模块的独立贡献。

\begin{figure}[htbp]
    \centering
    \resizebox{0.98\textwidth}{!}{%---------------------------------------------------------------------------%
%->> DP vs DCP decode latency comparison, DeepSeek-V2-Lite
%---------------------------------------------------------------------------%
\providecolor{benchAxis}{HTML}{222222}
\providecolor{benchGrid}{HTML}{D8DDE3}
\providecolor{benchDP}{HTML}{2F3A40}
\providecolor{benchDCP}{HTML}{8794A3}

\begin{tikzpicture}[
    x=1cm,
    y=1cm,
    axis/.style={benchAxis,line width=0.55pt},
    grid/.style={benchGrid,line width=0.35pt},
    dp/.style={benchDP,line width=0.9pt},
    dcp/.style={benchDCP,line width=0.9pt},
    tick/.style={font=\scriptsize,text=benchAxis},
    label/.style={font=\scriptsize,text=benchAxis},
    legend/.style={font=\scriptsize,text=benchAxis},
    panel/.style={font=\footnotesize,text=benchAxis},
]
\def\xA{0.95}
\def\xB{4.85}
\def\xC{8.75}
\def\xD{2.90}
\def\xE{6.80}
\def\ytop{4.45}
\def\ybottom{0.55}
\def\xstep{0.62}
\def\yscale{0.036}
\def\ymin{0}
\def\ymax{65}
\pgfmathsetmacro{\yheight}{(\ymax-\ymin)*\yscale}

\newcommand{\pt}[4]{({#1+#3*\xstep},{#2+(#4-\ymin)*\yscale})}
\newcommand{\drawaxes}[3]{
    \pgfmathsetmacro{\xright}{#1+4*\xstep}
    \node[panel,anchor=west] at (#1,#2+\yheight+0.22) {#3};
    \foreach \v in {0,20,40,60} {
        \pgfmathsetmacro{\ytick}{#2+(\v-\ymin)*\yscale}
        \draw[grid] (#1,\ytick) -- (\xright,\ytick);
        \draw[axis] (#1-0.05,\ytick) -- (#1,\ytick);
        \node[tick,anchor=east] at (#1-0.10,\ytick) {\v};
    }
    \foreach \i/\r in {0/1,1/2,2/4,3/8,4/16} {
        \pgfmathsetmacro{\xtick}{#1+\i*\xstep}
        \draw[axis] (\xtick,#2) -- (\xtick,#2-0.05);
        \node[tick,anchor=north] at (\xtick,#2-0.10) {\r};
    }
    \draw[axis] (#1,#2) -- (\xright,#2);
    \draw[axis] (#1,#2) -- (#1,#2+\yheight);
    \node[label,anchor=north] at (#1+2*\xstep,#2-0.34) {Request rate};
    \node[label,anchor=south,rotate=90] at (#1-0.50,#2+\yheight/2) {P90 latency (ms)};
}
\newcommand{\linefive}[8]{
    \draw[#3]
        \pt{#1}{#2}{0}{#4} --
        \pt{#1}{#2}{1}{#5} --
        \pt{#1}{#2}{2}{#6} --
        \pt{#1}{#2}{3}{#7} --
        \pt{#1}{#2}{4}{#8};
}
\newcommand{\solidmarks}[4]{
    \foreach \i/\v in {#4} {
        \fill[#3] \pt{#1}{#2}{\i}{\v} circle[radius=1.08pt];
    }
}
\newcommand{\openmarks}[4]{
    \foreach \i/\v in {#4} {
        \draw[#3,fill=white,line width=0.55pt] \pt{#1}{#2}{\i}{\v} circle[radius=1.08pt];
    }
}

\path[use as bounding box] (0.00,-0.08) rectangle (12.10,7.70);

% Legend
\draw[dp] (1.10,7.44) -- (1.58,7.44);
\fill[benchDP] (1.34,7.44) circle[radius=1.08pt];
\node[legend,anchor=west] at (1.72,7.44) {DP};
\draw[dcp] (2.55,7.44) -- (3.03,7.44);
\fill[benchDCP] (2.79,7.44) circle[radius=1.08pt];
\node[legend,anchor=west] at (3.17,7.44) {DCP};
\draw[axis,line width=0.9pt] (5.10,7.44) -- (5.58,7.44);
\fill[benchAxis] (5.34,7.44) circle[radius=1.08pt];
\node[legend,anchor=west] at (5.72,7.44) {TPOT};

\drawaxes{\xA}{\ytop}{4K baseline}
\drawaxes{\xB}{\ytop}{1\% 128K mix}
\drawaxes{\xC}{\ytop}{1\% 256K mix}
\drawaxes{\xD}{\ybottom}{1\% 512K mix}
\drawaxes{\xE}{\ybottom}{1\% 1M mix}

% 4K
\linefive{\xA}{\ytop}{dp}{6.58}{6.93}{6.98}{7.47}{8.55}
\linefive{\xA}{\ytop}{dcp}{7.15}{8.10}{9.00}{11.51}{22.88}
\solidmarks{\xA}{\ytop}{benchDP}{0/6.58,1/6.93,2/6.98,3/7.47,4/8.55}
\solidmarks{\xA}{\ytop}{benchDCP}{0/7.15,1/8.10,2/9.00,3/11.51,4/22.88}

% 128K
\linefive{\xB}{\ytop}{dp}{7.87}{8.73}{9.38}{10.13}{11.45}
\linefive{\xB}{\ytop}{dcp}{7.55}{8.33}{9.68}{12.53}{27.15}
\solidmarks{\xB}{\ytop}{benchDP}{0/7.87,1/8.73,2/9.38,3/10.13,4/11.45}
\solidmarks{\xB}{\ytop}{benchDCP}{0/7.55,1/8.33,2/9.68,3/12.53,4/27.15}

% 256K
\linefive{\xC}{\ytop}{dp}{8.97}{10.67}{11.94}{12.79}{15.72}
\linefive{\xC}{\ytop}{dcp}{7.56}{8.45}{9.44}{12.12}{29.33}
\solidmarks{\xC}{\ytop}{benchDP}{0/8.97,1/10.67,2/11.94,3/12.79,4/15.72}
\solidmarks{\xC}{\ytop}{benchDCP}{0/7.56,1/8.45,2/9.44,3/12.12,4/29.33}

% 512K
\linefive{\xD}{\ybottom}{dp}{13.39}{16.64}{18.24}{18.90}{22.20}
\linefive{\xD}{\ybottom}{dcp}{8.01}{9.15}{10.18}{13.40}{39.70}
\solidmarks{\xD}{\ybottom}{benchDP}{0/13.39,1/16.64,2/18.24,3/18.90,4/22.20}
\solidmarks{\xD}{\ybottom}{benchDCP}{0/8.01,1/9.15,2/10.18,3/13.40,4/39.70}

% 1M
\linefive{\xE}{\ybottom}{dp}{24.47}{28.09}{28.85}{30.25}{36.69}
\linefive{\xE}{\ybottom}{dcp}{9.19}{10.36}{11.96}{16.76}{60.89}
\solidmarks{\xE}{\ybottom}{benchDP}{0/24.47,1/28.09,2/28.85,3/30.25,4/36.69}
\solidmarks{\xE}{\ybottom}{benchDCP}{0/9.19,1/10.36,2/11.96,3/16.76,4/60.89}

\end{tikzpicture}
%---------------------------------------------------------------------------%
}
    \caption{DeepSeek-V2-Lite 长文本 Decode 场景下 DP 与静态 DCP 的 P90 延迟对比}
    \label{fig:dp-dcp-long-context}
\end{figure}

\textbf{正确性保证。}动态 CP 路径的推理输出与单 DP 路径在数值精度范围内一致，验证 CP attention 后端的 KV 分片与 LSE 聚合、跨 rank KV 虚拟化映射、TP 域去冗余、PD 传输重建和批次重排均不引入正确性偏差。

\textbf{平台适配验证。}在昇腾 910 系列集群上完成端到端部署和性能评估，量化昇腾平台硬件特征（HCCS 均匀带宽、高计算访存比）对动态 CP 收益边界和阈值选取的影响。

\section{拟采用的研究方法、技术路线、实验方案及其可行性分析}

\subsection{技术路线}

研究方法按"性能建模—系统原型—实验验证"三阶段推进。首先通过 profiling 建立"长请求 KV 访存 → DP rank 到达时间差 → MoE bubble"的定量瓶颈模型，明确动态 CP 的优化空间；其次按"单模块实现与验证 → 模块集成 → 端到端系统评估"的增量路径构建原型，每个阶段以输出正确性测试作为推进门槛；最后通过多基线对比实验和消融实验量化各模块的独立贡献与协同效果。

实现层面，改动主要涉及 vLLM 和 vLLM-Ascend 的四个层次：调度层（CrossDPScheduler、RequestQueue、负载均衡器）、KV Cache 层（CrossDPKVCacheManager、BlockTable）、执行层（EngineCoreProc、MultiProcExecutor、ModelRunner）和算子后端层（MLA CP attention、PCP 工具函数、AclGraph 图捕获）。其中调度层和 KV Cache 层的改动集中在 vLLM 侧，算子后端层的适配集中在 vLLM-Ascend 侧，执行层两端均需修改。

图~\ref{fig:dynamic-cp-system-architecture} 展示了动态 CP 系统架构。API Server 将请求转发至 DomainEngineCoreProc，其中 CrossDPScheduler 负责请求级 CP/DP 路径判定与跨 DP 协同调度，CrossDPKVCacheManager 完成跨 rank 的 KV Cache 虚拟化分配；调度结果以 list[SchedulerOutput] 的形式下发至 Dynamic CP Domain，由各 DP Worker 在统一的混合批次中执行 attention 计算与跨 rank 聚合。

\begin{figure}[htbp]
    \centering
    \includegraphics[width=\textwidth]{fig-dynamic-cp-system-architecture}
    \caption{动态 CP 系统架构}
    \label{fig:dynamic-cp-system-architecture}
\end{figure}
\FloatBarrier

技术路线按"瓶颈量化 → 动态判定与中心化调度 → 跨 DP KV 管理 → 混合批次执行 → 系统评估"五个阶段推进。

\begin{enumerate}
    \item \textbf{瓶颈量化阶段}：在昇腾环境中部署典型 MoE 长文本模型，测量 DP 与静态 CP 基线下的 TTFT、TPOT、ITL、QPS、HBM 读写量、AI Core 利用率、KV 使用量分布和 MoE bubble 时间，建立"长请求 KV 访存 → DP rank 到达时间差 → MoE bubble"的量化链路，并形成后续动态 CP 原型的对比指标。
    \item \textbf{动态判定与中心化调度阶段}：实现基于阈值 $T$ 的请求级 CP/DP 路径判定；实现 CentralizedScheduler，根据各 rank 剩余调度容量、KV 使用量和剩余 block 进行 rank 分配，统一生成各 rank 调度输出；搭建 Domain 执行架构（DomainEngineCoreProc），完成调度输出分发与跨 rank 元数据同步；解耦 Decode 阶段的 TP 与 CP，实现 CP 通信组的动态组建。
    \item \textbf{跨 DP KV 管理阶段}：实现 CrossDPKVCacheManager、二维 slot mapping、virtual block 和 mask 机制，验证 CP 请求跨 rank KV 分布的正确性；适配 TP 域内 KV 去冗余使其支持 CP/DP 请求共存；在 Mooncake connector 上扩展 PD 分离场景下的 DyCP KV 分片传输与 Decode 端重建。
    \item \textbf{混合批次执行阶段}：实现面向 MLA 的 CP attention 后端（mla\_cp），支持 latent KV 分片与 LSE 聚合；实现 CP 请求前置、DP 请求后置的批次重排与 hidden state restore；在昇腾图模式下探索二维 AclGraph 捕获与缓存策略。
    \item \textbf{系统评估阶段}：开展长短请求混合负载、不同阈值、不同 CP 请求比例、不同请求到达速率和不同通信拓扑下的实验，量化动态 CP 对 MoE bubble、尾延迟和吞吐的影响。
\end{enumerate}

\subsubsection{请求级动态 CP 判定与中心化跨 DP 调度}

\textbf{动态判定。}动态判定的核心是预设上下文长度阈值 $T$：对于新到达或进入调度周期的请求，调度器计算当前上下文长度 $L$。当 $L<T$ 时，请求以 CP=1 的 DP 路径执行；当 $L\geq T$ 时，请求进入 CP 路径，由多个 DP rank 协同执行。路径确定后，调度器根据各 rank 的剩余 KV block 和调度容量分配具体的执行 rank。阈值 $T$ 由离线 profiling 得到，根据模型结构、设备数量、硬件计算访存比和服务 SLO 设置。昇腾 NPU 的计算访存比约为 GPU 的 1.5--2 倍，Decode 阶段更易进入访存瓶颈区间，$T$ 可设得较低——较短的上下文即可从 CP 分摊获益。队列深度、KV 使用量和通信状态作为辅助约束，用于短请求 rank 选择、长请求并发控制和资源可用性检查。

\textbf{中心化调度。}为使 CP 请求跨 rank 同步执行，本文设计中心化跨 DP 协同调度替代传统"每个 rank 本地独立决策"的分布式调度。中心调度器维护每个 rank 的剩余调度容量（以可用 KV block 数为主，结合待处理请求数和队列深度）、KV 使用量和当前批次进度。每个调度周期中，长短请求分离队列分别管理两类请求：短请求按剩余容量和负载均衡选择单个 rank；长请求在满足并发度约束的前提下选择参与 CP 的 rank 集合，并保证同一 CP 请求在所有参与 rank 的调度周期中同步出现。CP rank 集合的选择结合通信拓扑：昇腾超节点架构中节点间与节点内通信带宽一致，CP group 可跨节点组建而不损失通信效率，支持更大 CP size 以覆盖超长上下文。调度结果包含请求路由、请求顺序映射、CP 区域标记和跨 rank 元数据同步信息。

\textbf{动态 CP 对 MoE 同步等待的缓解机制。}动态 CP 调度主要通过两类机制缓解 MoE bubble。一方面，短请求结合剩余 KV block 和 KV 使用量被路由到负载较轻的 rank，减少常规 DP 路径上的负载倾斜。另一方面，长请求不再压到单个 rank，而是通过 CP 路径把 attention 访存和计算分摊到多个 rank，缩小各 rank 到达 MoE all-to-all 的时间差。前期 DP 与静态 DCP 基线实验表明，低负载长上下文下 DCP 对尾延迟有收益，但高负载时静态 DCP 会因所有请求都通信而劣化。因此，动态判定的重点是只让超过阈值且资源允许的长请求进入 CP 路径。

\textbf{方案选择依据。}动态判定采用固定阈值而非自适应学习策略，原因在于 CP 的通信开销在两个阶段均可近似为常数：Prefill 阶段，需要切分的长请求长度通常远大于 chunk size，每步实际参与计算和通信的序列长度由 chunk size 决定，而 chunk size 在服务启动时已固定，因此 Prefill 阶段的 CP 通信量近似恒定；Decode 阶段，跨 rank 通信量很小，开销主要来自算子启动的固定延迟。两个阶段的通信开销均不随请求长度动态变化，只需通过离线 profiling 找到开启 CP 后能获得净收益的上下文长度拐点，即可在服务启动前确定阈值 $T$。第一版原型采用固定阈值以快速打通端到端链路，后续可基于实验数据进一步探索负载自适应调整。调度采用中心化架构而非分布式协商，原因在于 CP 请求必须在全部参与 rank 的同一调度周期同步出现，这一一致性约束是硬约束——分布式协商在异步环境下难以在有限通信轮次内保证多 rank 同步，而中心化调度器可以在单轮决策中生成全局一致的调度输出。

\textbf{Domain 执行架构。}中心化调度器产出的调度结果需分发至各 DP Worker 并协调跨 rank 执行，这一桥接由 Domain 执行架构（DomainEngineCoreProc）完成。Domain 是动态 CP 系统的多进程执行入口，一个 Domain 管理一组 DP Worker（通过 MultiProcExecutor），接收中心化调度器产出的 list[SchedulerOutput]，将各 rank 的调度输出分发至对应 Worker，并负责 CP 请求的跨 rank 元数据同步：确保参与同一 CP 请求的所有 rank 在同一执行步中获得一致的请求列表、CP 区域标记和通信域信息。将 Domain 作为调度与执行间的独立层，使 CrossDPScheduler 的输出格式与下游多进程执行器解耦，便于不同执行后端（Ascend/GPU）复用同一调度逻辑。此外，Domain 层承担 Decode 阶段 TP 与 CP 的解耦：现有 DCP 将 TP 和 CP 绑定，每个 rank 同时承担张量切分和序列切分职责，无法单独为不走 CP 的请求提供 TP-only 执行路径。本文设计将 TP 组和 CP 组分离管理，TP 组在服务启动时静态建立，CP 组由 Domain 根据每个调度周期中 CP 请求的参与 rank 集合动态组建。解耦后，同一组 DP Worker 既可以为 DP 请求执行仅含 TP 通信的 attention，也可以为 CP 请求在 TP 通信之上叠加 CP 通信，是混合批次执行的架构前提。

\begin{figure}[htbp]
    \centering
    \includegraphics[width=\textwidth]{fig-cross-dp-scheduler-flow}
    \caption{CrossDPScheduler 的长短请求分流与 rank 分配}
    \label{fig:cross-dp-scheduler-flow}
\end{figure}
\FloatBarrier

\subsubsection{跨 DP KV Cache 管理与二维映射}

动态 CP 的执行前提是同一请求的 KV Cache 可分布到多个 DP rank，而 attention 后端仍能按逻辑连续的 KV 序列访问。跨 DP KV Cache 虚拟化管理方案的设计要点是：每个 rank 保留独立的本地 KV block pool，与 vLLM 基于 PagedAttention 的内存管理方式兼容；对于 CP 请求，管理模块在逻辑层将多个 rank 的物理 block 组织为跨 DP 虚拟 KV 空间。核心机制为二维 slot mapping：第一维对应参与 CP 的 rank 编号，第二维对应逻辑 block 或 token 区间，映射关系包括本地物理 slot 位置、block 生命周期状态和跨 rank 通信参与标记。

这种二维映射支持三类访问模式。第一，DP 请求只占用某一个 rank 的本地区域，attention 后端按传统本地 block table 访问。第二，CP 请求按连续区间、block 轮转或交错粒度分布到多个 rank，attention 后端通过 virtual block 获取逻辑连续视图，再由底层映射定位到参与 rank 的物理分片。第三，同一批次中 CP 请求与 DP 请求共存时，mask 指示每个批次位置是否需要跨 rank attention 聚合，使通信算子只处理 CP 区域。

该方案面向动态 CP 所需的跨 DP KV 组织和访问问题，使 CP 请求和 DP 请求在同一批次中共享统一的缓存管理视图。采用二维 slot mapping 加虚拟 block 的方案而非全局共享内存池，是因为保留各 rank 独立的 block pool 可以兼容 vLLM 现有的 block 级内存管理逻辑，避免对底层内存分配器的侵入式修改；虚拟 block 仅在逻辑层提供跨 rank 的连续视图，物理层面各 rank 仍独立管理本地 KV block 的分配与回收。

\textbf{TP 域内 KV Cache 去冗余适配。}在原 DCP 实现中，同一 TP 组内的多个 rank 在 KV head 相同时只保留一份 KV 副本，通过 TP 通信在 attention 计算时共享，从而减少 TP 域内的 KV Cache 冗余。DyCP 引入 CP/DP 请求共存后，同一 TP 组内可能同时存在 CP 请求和 DP 请求：CP 请求的 KV 已按序列维度跨 DP rank 分片存储，每个 rank 只持有部分 KV；DP 请求的 KV 则完整保留在本地 rank。适配要点在于去冗余逻辑需按请求类型区分——DP 请求延续原有 TP 域去冗余策略，CP 请求的 TP 域去冗余需在跨 DP 分片后的局部 KV 上执行，确保不会误删 DP 请求的本地完整 KV 或遗漏 CP 请求分片 KV 的冗余副本。

\textbf{PD 分离场景下的 DyCP KV 传输适配。}在 PD 分离架构（Prefill 和 Decode 分别部署在不同实例）中，Prefill 实例生成的 KV Cache 需通过网络传输至 Decode 实例。DyCP 下，Prefill 产出的长请求 KV 按 CP 分片分布在多个 rank，与传统 PD 分离"每个请求 KV 完整存储在单 rank"的假设不同。方案在已有 Mooncake connector 基础上扩展按分片粒度的传输：每个 Prefill rank 仅传输本地持有的 KV 分片至 Decode 端对应 rank，Decode 端接收后通过 CrossDPKVCacheManager 的 slot mapping 重建跨 DP 虚拟 KV 视图，复用已有连接管理和序列化机制。

\FloatBarrier

\subsubsection{CP+DP 混合执行}

动态 CP 的技术挑战不仅是选择哪些请求走 CP 路径，还包括如何让 CP 请求和 DP 请求在同一个批次内高效执行。若 CP/DP 分别维护独立服务实例，会导致 NPU 资源闲置；若同一批次内所有请求统一执行 CP 通信，短请求延迟会恶化。这要求在 attention 计算适配、批次管理和图捕获三个层面分别处理。

\textbf{CP attention 后端适配。}不同注意力结构对 CP 分片和聚合的假设不同，CP attention 后端需要针对具体结构进行适配。当前首先适配的是 MLA（Multi-head Latent Attention）结构：MLA 将 KV 压缩为低维 latent 向量（compressed\_kv），每个 token 的 KV 存储量显著小于标准 MHA/GQA。CP 在 MLA 模型上执行时，分片的对象是 latent KV 而非标准多头 KV，分片粒度和通信内容需要适配——跨 rank 通信的是 latent attention 输出和 LSE，而非完整 KV 向量。CP attention 后端（mla\_cp）需要处理 latent 维度下的 KV 分片索引、attention mask 构建和跨 rank LSE 聚合。此外，chunked prefill 下同一批次中可能同时存在 Prefill CP、Prefill DP、Decode CP 和 Decode DP 四类请求，形成 2$\times$2 执行矩阵，需分别调用对应 CP 或标准 attention 后端并正确拼接输出。后续可按相同模式扩展至 GQA、MHA 等其他注意力结构。

\textbf{批次重排与 hidden state restore。}混合批次执行采用批次重排与通信区域限定机制：每个调度周期中，CP 请求排在批次前部，DP 请求排在批次后部，同时记录原始请求顺序与执行顺序的映射。attention 后端据 CP 区域边界标记只对批次前部执行跨 rank 通信。

Decode 阶段，CP 请求的每个参与 rank 只读取本地 KV 分片并计算局部 attention 输出 $O_r$ 和局部指数归一化统计量 $\operatorname{LSE}_r$。各 rank 交换或归约 $\operatorname{LSE}_r$，以数值稳定方式计算权重
\[
w_r=\exp(\operatorname{LSE}_r-\max_j \operatorname{LSE}_j),
\]
再聚合
\[
O=\frac{\sum_r w_r O_r}{\sum_r w_r},
\]
得到与未分片 attention 等价的全局输出。DP 请求则不参与该通信区域，仍在本地 rank 上按传统路径完成 attention，由此实现"短请求不通信，长请求才通信"。聚合完成后，需要将 hidden state 从执行顺序（CP 前置、DP 后置）恢复为原始请求顺序，以保证后续 FFN/MoE 层的输入与请求一一对应。

Prefill 阶段，将 CP 请求和 DP 请求的 query length、position id、block table、slot mapping 等元数据拆分，分别调用 CP attention 和 DP attention 后端，再按原始顺序合并输出。

\textbf{二维 AclGraph 捕获与缓存。}现有推理系统的图捕获通常是一维的，仅按 batch token 数或 batch size 进行 bucketing。动态 CP 引入 CP+DP 混合批次后，图的参数形状由 batch token 数和 CP 请求数两个独立维度共同决定：CP 请求数的变化会改变通信区域边界、跨 rank attention 聚合的参与范围和输出拼接方式，同一 batch token 数下不同 CP 请求数对应不同的计算图。若不加控制，二维组合空间会导致图缓存命中率下降和编译开销增大。方案设计面向 CP+DP 混合批次的二维图捕获与缓存策略，通过维度 bucketing 和常见组合预捕获，在图缓存大小与命中率之间取得平衡，使动态 CP 在图执行模式下可实际部署。

采用批次重排加通信区域限定而非将 CP/DP 请求分离到不同服务实例，是因为分离实例会导致各实例的 NPU 负载不均且资源利用率下降，而批次重排在同一批次内通过请求排列和 mask 标记实现通信隔离，不增加额外的实例管理和跨实例迁移开销。

\FloatBarrier

\subsection{实验方案}

实验模型以 DeepSeek-V2-Lite（16B/2.4B MoE，MLA 注意力）为主，该模型参数量适中且具备 MoE 结构，适合在有限卡数下验证动态 CP 对 MoE bubble 的缓解效果；条件允许时辅以 Qwen 等 GQA/MHA 模型验证方案泛化性。实验平台为 16 卡 Ascend 910C 集群（Atlas A3），节点内 HCCS 互联，使用 CANN、torch-npu、vLLM-Ascend 软件栈。实验负载覆盖纯短请求、纯长请求和长短混合三类场景，上下文长度覆盖 4K--512K（1M 视显存余量尝试），长请求比例设置为 0.1\%--10\%（对应真实服务中长请求占比较低的典型分布），请求到达速率覆盖 1--16 rps 区间。

围绕三个研究内容，实验分为以下三组。

\textbf{实验一：动态判定、中心化调度与 Domain 执行架构的有效性验证。}对比基线包括纯 DP（无 CP）、全局静态 CP、传统分布式本地 scheduler + min(req num) 路由、中心化但仅按请求数路由和 min(KV usage) 负载均衡方案。在不同上下文长度阈值 $T$、不同长请求比例和不同请求到达速率下，重点观察 MoE bubble 时间、ITL P90/P99 尾延迟和归一化吞吐的变化，验证动态判定相比静态配置的调度质量提升和 MoE 同步等待缓解效果；同时测量 Domain 执行架构的调度输出分发延迟和跨 rank 元数据同步开销。

\textbf{实验二：KV Cache 管理的正确性与开销验证。}包含三个子验证：（a）跨 DP KV 虚拟化——在 CP/DP 请求共存场景下，验证跨 rank KV 分片后的推理输出与单 DP 路径在数值精度内一致，测量二维 slot mapping 和 virtual block 的 KV block 利用率、显存峰值和分配回收开销，与 vLLM 默认的单 rank block table 管理基线对比；（b）TP 域 KV 去冗余在 DyCP 下的适配——验证 CP/DP 请求共存时 TP 域去冗余的正确性，量化 DyCP 下相比不去冗余的显存节省；（c）PD 分离场景下 DyCP KV 传输——测量按分片粒度传输的延迟开销，验证 Decode 端 KV 重建后的输出正确性。

\textbf{实验三：混合批次执行验证。}包含三个方面：（a）CP attention 后端——以 MLA 结构为首要验证对象，在不同 CP size 下验证 latent KV 分片与 LSE 聚合的数值精度，测量 CP attention 后端与标准 attention 后端的性能差异；（b）2$\times$2 执行矩阵——在 chunked prefill 开启时验证 Prefill/Decode $\times$ CP/DP 四类请求的元数据拆分与输出拼接正确性；（c）批次重排与图捕获——对比有/无批次重排对短请求尾延迟的影响，在不同 CP 请求数下测量二维图缓存命中率和编译开销，与一维 bucketing 基线对比。

\textbf{消融实验。}针对可替换的策略模块进行消融：（1）将动态判定替换为固定 CP size，观察静态配置下的性能退化；（2）将中心化调度替换为各 rank 本地独立调度，验证跨 DP 协同的必要性；（3）去除批次重排，使 CP/DP 请求不分区混合执行，观察短请求尾延迟变化。跨 DP KV 管理和通信约束属于正确性保障，不做去除消融。额外对比包括不同阈值 $T$ 和不同 CP rank 数量下的性能变化。

评价指标包括 TTFT、TPOT、ITL P50/P90/P95/P99、归一化 QPS、显存峰值、KV block 利用率、HBM 读写量及带宽利用率、AI Core 利用率、MoE bubble 时间、all-to-all 等待时间、CP attention 聚合通信时间、图捕获缓存命中率、请求顺序恢复正确性和推理输出一致性。

\subsection{可行性分析}

\textbf{理论可行性。}动态 CP 的核心计算——分布式 attention 聚合——基于在线 softmax 的数学等价性，该等价性已由 Flash-Decoding\citep{flashdecoding2023}、Ring Attention\citep{liu2023ringattention} 等工作阐述并在生产系统中广泛应用。各 CP rank 独立计算局部 attention 后，通过 LSE 加权聚合即可得到与单 rank 计算数值等价的结果。请求级动态判定仅改变请求的路由路径（CP 或 DP），不修改模型前向计算逻辑，因此不影响推理输出的正确性。

\textbf{技术可行性。}关键框架能力已在现有代码中可用：vLLM 的 Scheduler 和 BlockSpaceManager 提供了请求调度和 block 级 KV 管理接口，vLLM-Ascend 已集成 PCP/DCP attention 后端（mla\_cp、pcp\_utils）和 HCCL 通信组管理，AclGraph 模块支持昇腾图模式的捕获与缓存。对于新增子点：TP 域 KV 去冗余已在 vLLM-Ascend 的 DCP 实现中完成，DyCP 适配为增量修改；Mooncake connector 已在 vLLM-Ascend 中实现 PD 分离的 KV 传输协议，DyCP 传输适配可复用其连接管理和序列化机制；MLA 下的 CP attention 后端（mla\_cp）已具备基本实现，DyCP 适配需在此基础上扩展混合批次支持。前期已在 vLLM 社区提交 Dynamic Context Parallelism RFC\citep{vllmdynamiccp2025}，完成了 CrossDPScheduler、CrossDPKVCacheManager 和 Domain 等核心模块的接口设计。前期在 16 卡 Ascend 910C 上完成的 DP 与静态 DCP 对比实验（图~\ref{fig:dp-dcp-long-context}）已验证了 CP 在长上下文场景下的性能收益和静态配置的局限性，为动态 CP 的阈值选取提供了实验依据。

\textbf{条件可行性。}实验室具备 16 卡 Ascend 910C 集群（Atlas A3），节点内 HCCS 互联，已在该集群上完成 DeepSeek-V2-Lite 的 DP 和静态 DCP 基线实验。软件栈方面，CANN、torch-npu、vLLM-Ascend 已完成部署，具备自定义算子编译和 AclGraph 图模式执行的开发调试环境。

\textbf{主要风险与应对措施。}（1）跨 rank 协同调度对一致性要求高，若某个 rank 的批次缺失 CP 请求，会导致通信死锁或结果错误；对此通过中心化调度器统一生成所有 rank 输出，并在执行前校验 CP 请求集合和通信域。（2）跨 DP KV mapping 容易因 block 状态、slot 偏移或顺序恢复错误导致输出异常；对此采用先单请求、再多请求、再混合批次的渐进验证，并用输出一致性测试对比单 DP 路径。（3）CP 通信可能在高负载下成为瓶颈；对此保留 DP 路径作为退路，在通信拥塞或长请求比例过高时减少 CP 请求调度量，以静态全局 CP 的通信量作为限流上界。

\section{已有科研基础与所需的科研条件}

\subsection{已有科研基础}

前期调研已覆盖阶段异构调度、序列并行、长上下文推理和 vLLM-Ascend CP 工程实现等方向，相关工作包括 Splitwise、DistServe、Mooncake、Ring Attention、Ulysses、USP、百万 token 推理 CP 等，初步明确了 CP 并行、KV Cache 管理、跨 DP 调度和 MoE 同步等待之间的关系。作者已在 vLLM 社区提交 Dynamic Context Parallelism RFC\citep{vllmdynamiccp2025}，完成了动态 CP 的方案设计；本文将在该设计基础上完成系统实现和实验验证。

工程方面，已阅读并分析 vLLM-Ascend 中 PCP/DCP 的实现代码，包括 rank layout 构建、slot mapping 计算、attention 后端调用和 HCCL 通信组初始化等路径，理解 PCP/DCP 与 DP、TP 组合时 block table、position id 和 hidden state 的元数据变化。前期还围绕 PD 分离场景下的 KV Cache 布局与重排进行过工程验证，覆盖 TP 异构映射、PP/MTP 层划分等问题。这部分不作为独立研究目标，但相关经验可直接复用于 CrossDPKVCacheManager 的 slot mapping 和 block 状态管理实现。

实验方面，已在 16 卡 Ascend 910C 集群上以 DeepSeek-V2-Lite 完成 DP 与静态 DCP 的基线对比（图~\ref{fig:dp-dcp-long-context}），覆盖 4K、128K、256K、512K、1M 上下文和 1--16 rps 负载。结果表明，低负载下 DCP 可降低长上下文的 P90 尾延迟，但高负载下静态 DCP 因短请求也参与通信而劣于纯 DP。这一交叉现象直接驱动了本文的核心设计：请求级动态选择 CP/DP 路径。

\subsection{所需科研条件}

所需科研条件主要包括硬件平台、软件环境、模型与性能分析工具。硬件方面，需要多张昇腾 NPU 的服务器或集群，并具备 HCCS/RDMA 等高速互联能力。软件方面，需要安装匹配版本的 CANN、驱动固件、torch-npu、vLLM-Ascend、MindIE 及相关推理框架，并具备自定义算子编译和图模式执行的开发调试环境。模型方面，准备 DeepSeek、Qwen 等长上下文 MoE 模型及其量化版本。工具方面，需要日志、profiling、显存监控、通信监控和 benchmark 工具，采集 TTFT、TPOT、ITL、QPS、HBM 读写量、MoE bubble 和通信开销等指标。

\section{研究工作计划与进度安排}

\begin{enumerate}
    \item 2026 年 6 月至 2026 年 8 月：完成文献调研、系统调研和基础实验环境搭建。
    \begin{itemize}
        \item 梳理 CP 并行、MoE 推理、KV Cache 管理、跨 DP 调度和昇腾推理框架相关工作。
        \item 完成昇腾服务器环境部署，验证 vLLM-Ascend 和静态 CP 基本功能。
    \end{itemize}
    \item 2026 年 9 月至 2026 年 11 月：完成动态 CP 判定、中心化跨 DP 调度与 Domain 执行架构原型。
    \begin{itemize}
        \item 实现基于阈值 $T$ 的 CP/DP 路径判定。
        \item 实现 CentralizedScheduler 的全局状态维护、长短请求队列和按 rank 调度输出。
        \item 搭建 Domain 执行架构（DomainEngineCoreProc），完成调度输出分发与跨 rank 元数据同步；解耦 Decode 阶段的 TP 与 CP，实现 CP 通信组的动态组建。
    \end{itemize}
    \item 2026 年 12 月至 2027 年 1 月：完成跨 DP KV Cache 管理。
    \begin{itemize}
        \item 实现 CrossDPKVCacheManager、二维 slot mapping、virtual block、mask 和引用状态管理。
        \item 适配 TP 域内 KV 去冗余使其支持 CP/DP 请求共存场景。
        \item 在 Mooncake connector 上扩展 PD 分离场景下的 DyCP KV 分片传输与 Decode 端重建。
        \item 完成单请求 CP、多请求 CP、CP+DP 共存场景下的 KV 正确性和输出一致性验证。
    \end{itemize}
    \item 2027 年 2 月至 2027 年 3 月：完成 CP+DP 混合批次执行与昇腾图模式适配。
    \begin{itemize}
        \item 实现 CP attention 后端，首先适配 MLA 结构的 latent KV 分片与 LSE 聚合。
        \item 实现 CP 请求前置、DP 请求后置的批次重排、hidden state restore 和 Decode 通信区域限定。
        \item 实现二维 AclGraph 捕获与缓存，并补齐 chunked prefill 下 2$\times$2 执行矩阵的元数据拆分执行。
    \end{itemize}
    \item 2027 年 4 月：开展系统实验、消融实验和长短混合负载实验，整理论文主体和图表。
    \item 2027 年 5 月：完成论文修改、预答辩和正式答辩材料准备，完成毕业归档。
\end{enumerate}

%---------------------------------------------------------------------------%
